% --- s01.png ---
\underline{Complex projective varieties and quantum mechanics}

\underline{Definition:}\quad
\[
\mathbb{CP}^n=\left\{(z_0,\ldots,z_n)\ \middle|\ \sum_{j=0}^{n}|z_j|^2=1\right\}
\]

We associate $\mathbb{CP}^n$ to a qudit of dimension $n$.

The many body Hilbert space we work with is
\[
L^2(\mathbb{CP}^n)=\bigoplus_{N\geq 0}h_N
\]
which is graded by degree,
\[
\deg(p(x))=N
\]
of homogeneous polynomials.

\[
\lvert k_0,k_1,\ldots,k_n\rangle
=
\sqrt{\frac{(k_0+k_1+\cdots+k_n+n)!}{k_0!k_1!\cdots k_n!n!}}\,
z_0^{k_0}z_1^{k_1}\cdots z_n^{k_n}
\]

Forms orthonormal basis w.r.t.
\[
\langle k\mid l\rangle
=
\int_{\mathbb{CP}^n}\overline{\underline{z}}^{\,\underline{k}}\underline{z}^{\underline{l}}\,d\underline{z}
\]

An ideal (graded) is comprised of all homogeneous polynomials $p(z)$ s.t.
\[
(f_1,\ldots,f_d)_N
=
\left\{
p(z)=f_1(z)h_1(z)+\cdots+f_d(z)h_d(z)
\ \middle|\ 
\deg p=N
\right\}
\]

We want to realize $(f_1,\ldots,f_d)_N$ as ker of Hamiltonian. We need introduce (unnormalised) creation \& annihilation operators:
\[
a_j^\dagger=z_j,\qquad a_j=\partial_j
\]

Also write
\[
\ell_j^\dagger\lvert k_0,\ldots,k_n\rangle
=
\lvert k_0,\ldots,k_j+1,\ldots,k_n\rangle
\]
\[
\ell_j\lvert k_0,\ldots,k_n\rangle
=
\lvert k_0,\ldots,k_j-1,\ldots,k_n\rangle
\]

We have
\[
\ell_j\propto a_j,\qquad \ell_j^\dagger\propto a_j^\dagger
\]

% --- s02.png ---
The vector subspace $V_{(f_1,\ldots,f_d)_N}^{\perp}$ is defined by the property:
\[
q(z)\in V_{(f_1,\ldots,f_d)_N}^{\perp}
\Longleftrightarrow
\langle q(z)\mid p(z)\rangle=0
\quad \forall\,p(z)\in V_{(f_1,\ldots,f_d)_N}
\]

\[
p(z)\in V_{(f_1,\ldots,f_d)_N}
\Longrightarrow
\]

\[
p(z)=\sum_{j=1}^{d}h_j(z)f_j(z)
\]

\[
=\sum_{j=1}^{d}a^\dagger(f_j)\lvert h_j\rangle
\]

where
\[
a^\dagger(f)=\sum_{k_0,\ldots,k_n}f_{k_0\ldots k_n}
(a_0^\dagger)^{k_0}(a_1^\dagger)^{k_1}\cdots(a_n^\dagger)^{k_n}
\]

So
\[
\langle q(z)\mid p(z)\rangle=0
\Longrightarrow
\]

\[
\sum_{j=1}^{d}
\langle q(z)\mid a^\dagger(f_j)\lvert h_j\rangle
=0
\qquad \forall\,h_j(z)
\]

\[
\Longrightarrow
\sum_{j=1}^{d}
\langle q\mid a^\dagger(f_j)\lvert h_j\rangle
=0
\qquad \forall\,\lvert h_j\rangle
\]

\[
\Longleftrightarrow
a(f_j)\lvert q\rangle=0
\qquad \forall\,j=1,\ldots,d
\]

\[
\Longleftrightarrow
\lvert q\rangle\in\ker H
\quad\text{where}
\]

\[
H=\sum_{j=1}^{d}a^\dagger(f_j)a(f_j)
\]

Check with example

\[
f_0=z_0z_1-z_2^2
\]

\[
a(f_0)=a_0a_1-a_2^2
\]

\[
[a_0,a_0^\dagger]z_0^k
\]

\[
[\partial_0,z_0]:
\quad
\partial_0(z_0z_0^k)-z_0\partial_0z_0^k
\]

\[
=(k+1)z_0^k-kz_0^k
\]

\[
=z_0^k
\]

% --- s03.png ---
\[
(a_0a_1-a_2^2)(z_0z_1-z_2^2)
=1-2\mathrel{\overset{?}{=}}-1
\]

Note:

\[
a(f_0)\lvert 00\cdots\rangle=0
\]

\[
a(f_0)\lvert k_0\cdots k_n\rangle=0
\qquad\text{if }\sum_{j=2}^n k_j=1
\]

\[
a(f_0)\lvert 112\rangle
=\lvert 002\rangle-2\lvert 110\rangle
\]

\[
a(f_1)\lvert 11\rangle=\lvert 000\rangle
\]

\[
a(f_0)\lvert 002\rangle=-2\lvert 000\rangle
\]

\[
\Longrightarrow\quad
a(f_0)\lvert\psi\rangle=0
\quad\Longrightarrow\quad
\lvert\psi\rangle=\lvert 110\rangle+\frac{1}{2}\lvert 002\rangle
\]

But this is in $\mathcal H$? We realize $V^\perp_{(f_1,\ldots,f_d)_N}$ as $\ker H$.

Note $V_{(f_1,\ldots,f_d)_N}$.

Is it helpful to have a ground state space given by $V_I^\perp$? Consider, e.g., ideal membership problem:

Suppose we have $\ker H=V_I^\perp$. Is $f\in I$

Basic:
\[
H+L_f^\dagger L_f=H'
\]

If $f\in I$ then
\[
\ker H=\ker(H+H')
\]

If $f\notin I$ then
\[
\ker(H+H')=?
\]

% --- s04.png ---
If $f\notin I$ then
\[
\ker(H+H')=? 
\]

\[
\ker(H+H')=
\left\{\lvert g\rangle\ \middle|\ 
\left\langle g\middle|\sum_j h_jf_j+hf\right\rangle=0
\quad\forall h_1,\ldots,h_d,h
\right\}
\]

i.e. $f\notin I$ then
\[
I\subset (I,f)
\]
and $(I,f)$ is no bigger. So
\[
V_{(I,f)}\text{ is smaller.}
\]
That is
\[
\ker(H+H')\subset \ker(H).
\]

Perturbation theory would make value different in limit $p\to0$.

Perhaps it is worth trying to build a parent Hamiltonian with
\[
\ker(H)=V_I.
\]
This should allow finite field.

Strategy: First find parent for single ideal then grow the ideal and add constraints.

Hamiltonian for $(f)$?

\[
I=(f)=\{h(x)f(x)\mid h(x)\in k[z]\}
\]

We start with:
\[
H=\mathbb I-\lvert f\rangle\langle f\rvert
\]
\[
=\mathbb I-L_f^\dagger\lvert0\rangle\langle0\rvert L_f
\]

Would be good to have unitary version of $L_f^\dagger$.

\[
\ker H=\{\lvert g\rangle\mid \lvert\langle g\mid f\rangle\rvert=1\}
\]

Look now at
\[
H=\sum_{h(x)}\mathbb I-\lvert hf\times hf\rvert^2
\mathrel{\overset{?}{=}}
\sum_h\mathbb I-\lvert u_hf(x)\rangle\langle u_hf(x)\rvert
\]

\[
(1)=\mathbb R\quad\Longrightarrow\quad H_{(1)}=0
\]

% --- s05.png ---
But what is $u_f$? Can we use $L_f$ to do this?

For $|f\rangle$:

What $\ker H=\{\,|g\rangle\mid |\langle f|g\rangle|=1\,\}$

What equation is satisfied so that
\[
M_f|f\rangle=0
\]
\[
\Rightarrow L_f-1
\]

\[
f(x)=\sum_{j=0}^{\deg(f)} f_jx^j
\]

\[
|f\rangle=\sum_{j=0}^{\deg(f)} f_j|j\rangle
\]

Full basis for $\mathbb{CP}^n$:

\[
\mathcal{H}=\operatorname{span}\left\{|\alpha_0,\ldots,\alpha_n\rangle\mid \alpha_i\in\mathbb{Z}^+;\ |\alpha\rangle\leftrightarrow z^\alpha\right\}
\]

\[
|f\rangle\longmapsto |h_g\rangle\;?
\]

\[
f(x,y)=\sum_{j+k=d}f_{jk}x^jy^k
\]

\[
h(x,y)=\sum_{l,m}h_{lm}x^ly^m
\]

\[
f(x,y)h(x,y)=\sum_{\substack{j,k\\j+k=d}}\sum_{l,m}f_{jk}h_{lm}x^{j+l}y^{k+m}
\]

\[
|fh\rangle=\sum_{\substack{j,k\\j+k=d}}\sum_{l,m}f_{jk}h_{lm}|j+l,k+m\rangle
\]

If $h$ is homogeneous.

% --- s06.png ---
\[
|fh\rangle=\sum_{\substack{j,k\\j+k=d}}\sum_{l,m}f_{jk}h_{lm}|j+l,k+m\rangle
\]

If $h$ is homogeneous:

\[
|p\rangle \xleftarrow{h_{\deg(f)}} \xrightarrow{h} |hf\rangle\in\mathcal{P}_{\deg(h)+\deg(f)}
\]

First look at operator \quad $f\mapsto x^\ell y^m$

\[
|j,k\rangle\xrightarrow{x^\ell y^m}|j+\ell,k+m\rangle
\]

This is a shift-right operator. \quad $J^\ell|j\rangle=|j+\ell\rangle$

\[
J=\sum_{j=0}^{\infty}|j+\ell\rangle\langle j|
\]

So this makes Hilbert space interpretation:

\[
\mathcal{H}_R\equiv
\underbrace{\ell^2(\mathbb{Z}^+)\otimes\ell^2(\mathbb{Z}^+)\otimes\cdots\otimes\ell^2(\mathbb{Z}^+)_j}_{m+1}
\]

$J$ is not natural in a Harmonic oscillator setting.

\[
\ell^2(\mathbb{Z})
\]

would then be appropriate for Laurent polynomials.

So we mirror $\mathbb{F}[x]$ case:

We want to create superposition of all polys in $V_I$

\[
|b\rangle\longmapsto\sum_{h\in R_d}|hf\rangle
\qquad;\qquad
R=\bigoplus_{d=0}^{\infty}R_d
\]

Is arbitrary enough?

\[
|f\rangle\longmapsto|f+h\rangle\Rightarrow\infty
\]

% --- s07.png ---
Is addition unitary?

\[
|f\rangle\longmapsto |f+h\rangle\Rightarrow\infty
\]

But we could start with forming $|k\rangle$.

Thus

\[
R=\sum_{\alpha_0,\ldots,\alpha_n}
|\alpha_0,\alpha_1,\ldots,\alpha_n\rangle
\]

Assume this is easy to create (how?)

Remember! $|0\cdots 0\rangle$ stands for ``$1$'' etc.

\[
1\longmapsto z_0^{\alpha_0}z_1^{\alpha_1}\cdots z_n^{\alpha_n}
\]

\[
|0\cdots 0\rangle\longmapsto
L_{z_0}^{\alpha_0}\cdots L_{z_n}^{\alpha_n}|0\cdots 0\rangle
=|\alpha_0,\alpha_1,\ldots,\alpha_n\rangle
\]

$\therefore$ we could proceed by applying:

\[
\mathrm{CG}=\sum_{j=0}^{\infty}J^j\otimes |j\rangle\langle j|
\]

To form

\[
l(z^\alpha)=\sum_\beta|\alpha+\beta\rangle
\]

\[
=\sum_\beta J^\beta|\alpha\rangle
\]

\[
l(z_1^{\alpha_1}z_2^{\alpha_2})
=\sum |\alpha_1+\beta_1\rangle\oplus|\alpha_2+\beta_2\rangle
\]

Eg. $\alpha_1=(1,0)\quad \alpha_2=(0,1)$

We need to form superpositions:

\[
|0\rangle\longmapsto\frac{1}{\sqrt{2}}
\left(|\alpha_1\rangle+|\alpha_2\rangle\right)
\]

Actually: this is ok for monomial ideals:

\[
(z_1^{\alpha_1},z_2^{\alpha_2})
=\{p(z)\mid p(z)=h_1(z)z^{\alpha_1}+h_2(z)z^{\alpha_2}\}
\]

\[
|p(z)\rangle\in\operatorname{span}
\left\{
|\alpha_1+\beta_1\rangle,\,
|\alpha_2+\beta_2\rangle
\mid \beta_1,\beta_2\in(\mathbb{Z}^+)^{n+1}
\right\}
\]

% --- s08.png ---
\[
(z^{\alpha_1},z^{\alpha_2})
=
\left\{p(z)\mid p(z)=h_1(z)z^{\alpha_1}+h_2(z)z^{\alpha_2}\right\}
\]

\underline{Example.}\quad
\[
|(z_1,z_2)\rangle\in\operatorname{span}
\left\{
|1+\beta_1^0,\beta_2^0\rangle,\,
|\beta_1^0,1+\beta_2^0\rangle
\right\}
\]

\[
\in\mathcal{H}_N\cap\{|0,0\rangle\}^{\perp}
\]

The space $|f\cdot h\rangle$ is large enough to characterise

\[
|f\cdot h\rangle\in\operatorname{span}
\left\{
J^\beta\left(\sum_\alpha f_\alpha|\alpha\rangle\right)
\mid \beta\in(\mathbb{Z}^+)^{n+1}
\right\}
\]

So, really, we want to start with

\[
\sout{H=\mathbb{1}-\sum_{j=1}^{d}\lvert f_j\rangle\langle f_j\rvert}
\]

\[
\Longleftarrow \underline{\text{AND}}\text{ quantifier (!)}
\]

No, actually, OR quantifier

\[
H\lvert f_j\rangle=0
\]

Then sum over all shifts:

\[
H\longmapsto\sum_\alpha J^\alpha H(J^\alpha)^\dagger
\]

\[
H\sum_j c_j\lvert f_j\rangle=
\]

\[
\sum_j c_j H\lvert f_j\rangle=0
\]

(Can drop $\mathbb{1}$; No, normalisation not important.
